% Separate post-CP30 candidate; B85 dependency awaits its own promotion.
\subsection{The cyclotomic quotient excludes a traceless six-dimensional image}

Use the complete discretely valued field, actual torsion
comparison and tame good extension \(L/K\) of B85:
\(p\ge5\), \(e_K=p(p-1)\), algebraically closed residue,
\(\mu_p\subset K\), \([L:K]=N=p^2-1\).
All valuations below have value group \(\mathbf Z\) on
the field named in their subscript. Put \(q=p^2\).

\begin{lemma}[The last cyclotomic break after the tame good extension]
\label{lem:B-RD86-cyclotomic}
\status{Candidate; complete proof supplied, independent audits pending.}
If \(\mu_{p^2}\subset L\), then \(\mu_{p^2}\subset K\),
\([L(\mu_{p^3}):L]=p\), and every nonidentity element
of that cyclic group has lower break \(N^2\).
\end{lemma}
\begin{proof}
Because \(\mu_p\subset K\), the degree of
\(K(\mu_{p^2})/K\) divides \(p\). It also divides
the prime-to-\(p\) degree \(N\), so is one.
Choose a Cohen coefficient ring \(\Lambda\subset\mathcal O_K\)
and put \(K_0=\operatorname{Frac}(\Lambda)\);
its existence follows from the full Cohen structure
theorem \cite[Tag 032A]{Stacks032A}.
The residue field is perfect, so \(\Lambda\) is an
unramified complete discrete valuation ring with
uniformizer \(p\) and the same residue field.

For a primitive \(p^2\)-th root \(\zeta_2\), the
polynomial \(\Phi_{p^2}(1+T)\) is Eisenstein over
\(\Lambda\), of degree \(p(p-1)\). Its root
\(\zeta_2-1\) is therefore a uniformizer of \(K\).
The complete subring \(\Lambda[\zeta_2]\) contains
fixed representatives of every residue element and
that uniformizer. Successive uniformizer expansion
therefore gives \(\mathcal O_K=\Lambda[\zeta_2]\),
and \(K=K_0(\zeta_2)\). This explicitly justifies
the coefficient-field identification at the exact
absolute index; no canonical choice of \(\Lambda\)
is required.

Let \(K_3=K_0(\zeta_3)\), where \(\zeta_3\) has
order \(p^3\) and \(\zeta_3^p=\zeta_2\).
The analogous Eisenstein polynomial has degree
\(p^2(p-1)\), and \(K_3/K\) is cyclic of degree
\(p\). A uniformizer is \(\pi=\zeta_3-1\).
For a nonidentity automorphism \(\sigma\),
\[
   \sigma\pi-\pi=\zeta_3(\zeta_p^a-1),\quad
   p\nmid a,\qquad v_{K_3}(\sigma\pi-\pi)=p^2.
                                                        \tag{B-RD86.1}
\]
Thus its lower break is \(p^2-1=N\).

The extensions \(L/K\) and \(K_3/K\) are disjoint.
Their composite \(C=LK_3\) is the tame extension of
\(K_3\) of degree \(N\). Since the residue field is
algebraically closed, every unit has an \(N\)-th root
by Hensel's lemma. The tame extension consequently
admits a uniformizer \(t\) with \(t^N=\pi\);
this is also the direct tame uniformizer argument in
\cite[Proposition 2.12 and proof, p.~11]{ColmezCorpsLocaux}.
Extend \(\sigma\) by fixing \(L\). Since it is wild,
\(\sigma(t)/t\) has residue one. The identity
\[
       \left(\frac{\sigma(t)}t\right)^N
           =\frac{\sigma\pi}{\pi}
                                                        \tag{B-RD86.2}
\]
and the unit \(N\) show that
\[
  v_C\left(\frac{\sigma(t)}t-1\right)
   =v_C\left(\frac{\sigma\pi}{\pi}-1\right)
   =N\,v_{K_3}\left(\frac{\sigma\pi}{\pi}-1\right)
   =N^2.                                               \tag{B-RD86.3}
\]
This is the asserted lower break of \(C/L\).
\end{proof}

\begin{proposition}[The two actual congruence layers cannot both be traceless]
\label{prop:B-RD86-traceless}
\status{Candidate; complete proof supplied, independent audits pending.}
Suppose the residual representation is tame reducible,
the finite determinant is nonsquare, the high lines
are exchanged, and the actual parameters are
\(B_i=p^2k_i\), \(1\le k_i\le p-2\).
Assume the residual eigenvalue ratio is different
from one. Then the common actual third-level image
over \(L\) cannot have
\[
                  W_1=W_2=\mathfrak{sl}_2(\mathbf F_p)
                                                        \tag{B-RD86.4}
\]
in the basis \(H_1=\mathbf F_pe_1\),
\(H_2=\mathbf F_pe_2\).
\end{proposition}
\begin{proof}
Suppose (B-RD86.4) holds. On the good field \(L\)
the actual determinant is the cyclotomic character
\(\chi_p\), as in B76 and B77; the finite determinant
character is killed by the good extension. The first
traceless layer makes this determinant trivial modulo
\(p^2\), and therefore \(\mu_{p^2}\subset L\).
The lemma shows that the remaining character modulo
\(p^3\) is nonzero of order \(p\).

For every terminal element,
\[
       \det(1+p^2A)=1+p^2\operatorname{tr}(A)=1
                        \pmod{p^3}.                  \tag{B-RD86.5}
\]
Hence the actual determinant induces a nonzero linear
functional
\[
 \ell:W_1=\mathfrak{sl}_2(\mathbf F_p)\longrightarrow
 \mathbf F_p,\qquad
 \det(g)=1+p^2\ell(A)\pmod{p^3},\quad g\mapsto A.
                                                        \tag{B-RD86.6}
\]
This does not assert that the determinant itself is
zero merely because its terminal restriction is zero.
Its invariance under actual tame conjugation implies
\(\ell(E_{12})=\ell(E_{21})=0\), since the two
off-diagonal weights are different from one.
Writing \(H=\operatorname{diag}(1,-1)\), one has
\(\ell(H)\ne0\).

Choose an actual primitive second-level point \(P\)
with \(pP=e_1\), and let \(E=L(P)\). The common
second-level action is translation by \(W_1e_1\),
which is all of \(\mathbf F_p^2\). All first-level
points are rational, so \(E/L\) is Galois of degree
\(q\), with
\[
 \operatorname{Gal}(E/L)=
      \mathfrak{sl}_2(\mathbf F_p)/\mathbf F_pE_{12}.
                                                        \tag{B-RD86.7}
\]
The actual point stabilizer also contains the terminal
layer, on which the determinant vanishes. By (B-RD86.6)
the entire point stabilizer is contained in the
cyclotomic kernel. Thus the cyclotomic field
\(C=L(\mu_{p^3})\) is a subfield of this actual point
field \(E\).

Every class in (B-RD86.7) has a representative
\(aH+cE_{21}\). It translates \(P\) by
\(ae_1+ce_2\), and it has nontrivial image in
\(\operatorname{Gal}(C/L)\) exactly when \(a\ne0\).
In model 2 such a translation is outside its high
line \(\mathbf F_pe_2\), hence has low valuation
\(pk_2\) over \(L\). The target \(pP=e_1\) is
also low in model 2, so its inverse coordinate has
valuation \(k_2/p\) over \(L\). In \(E\),
\[
       v_E(T_2(P))=pk_2,\qquad
       v_E(gT_2(P)-T_2(P))=p^3k_2\quad(a\ne0).
                                                        \tag{B-RD86.8}
\]
The unconditional lower-break inequality of B60 gives
\[
             b_E(g)\le p^3k_2-pk_2=pNk_2
                \le pN(p-2)<N^2                   \tag{B-RD86.9}
\]
for every lift of every nonidentity element of
\(\operatorname{Gal}(C/L)\).

For completeness the exact quotient formula for a
nonidentity \(\bar g\in\operatorname{Gal}(C/L)\) is
\[
 b_C(\bar g)+1
   =\frac1{[E:C]}\sum_{g\mapsto\bar g}(b_E(g)+1).
                                                        \tag{B-RD86.10}
\]
Here both fields are totally ramified and the fibre
has \([E:C]=p\) elements. Formula (B-RD86.10) is the
norm-polynomial formula, with its full two-divisibility
proof, in \cite[Proposition 3.7, pp.~14--15]{ColmezCorpsLocaux}.
It uses no finite-residue hypothesis. It implies from
(B-RD86.9) that \(b_C(\bar g)<N^2\), contradicting
the cyclotomic lemma. This proves the proposition.
\end{proof}

\begin{corollary}[Seven dimensions in the exchanged tame sector]
\label{cor:B-RD86-tame}
\status{Candidate; complete proof supplied, independent audits pending.}
Assume \(p\equiv1\pmod4\). In the different-character
tame reducible sector at \(e=p(p-1)\), an actual
\(p^3\)-congruent pair must have wild image order at
least \(p^7\). At \(p=5,e=20\), the necessary
\(\{25,75\}\) tame parameter pattern can therefore
remain only at wild orders \(5^7\) or \(5^8\).
\end{corollary}
\begin{proof}
B85 gives \(\dim W_1\ge3\), and actual powers give
\(W_1\subset W_2\). If the total dimension were at
most six, both spaces would have dimension three and
would coincide. B83 puts \(E_{12},E_{21}\) in the
first space, and their actual commutator puts
\(H=\operatorname{diag}(1,-1)\) in the second.
The two equal three-dimensional spaces would therefore
both be \(\mathfrak{sl}_2(\mathbf F_p)\), excluded by
the proposition. The total dimension is at least seven;
the general two-layer upper bound is eight.
The parameter restriction is B84, with no existence
assertion.
\end{proof}

\begin{remark}[The remaining image orders and source scope]
\label{rem:B-RD86-scope}
\status{Candidate scope; no full first-wild least-index claim.}
This argument does not exclude the seven- or
eight-dimensional images and does not strengthen the
other two residual sectors. The full first-wild
least-index interval remains \(\{2,3\}\).
The new obstruction uses the actual cyclotomic
subfield of a degree-\(p^2\) point field. It is not a
claim that all traceless terminal images have zero
finite determinant.

The complete original Cohen structure proof, Tag 032A,
and Colmez's original Proposition 2.12 with proof on
p.~11 and Proposition 3.7 with its full proof on
pp.~14--15 were freshly opened and read for this
candidate. The elementary cyclotomic and tame-root
calculations are written above. The actual B83/B85
matrix layers and the full B60 displacement lemma
were reread. The separate B85 dependency must itself
have completed audits before this corollary is
promoted. All current proofs remain outside the
frozen checkpoint30.

\smallskip
\noindent\emph{Exact supplement.}
The retained standalone program
\path{code/reduction_descent/verify_first_wild_cyclotomic_quotient.py}
checks the symbolic determinant, bracket, point-field
evaluation and strict-gap identities, every small
parameter for primes 5 through 101, and all diagonal
characters on the displayed additive point quotients
for primes through 31. It also computes every
coefficient of the six shifted cyclotomic polynomials
at \(p=5,7,11\), \(n=2,3\), and verifies their
Eisenstein conditions. An explicit nonzero determinant
with traceless terminal power is retained as a guard
against the invalid simpler inference. The author run
passed in \(0.4382091660081642\) seconds. No execution failed.
The finite ranges do not replace the actual field
inclusion and quotient-ramification proofs. Every
executed program, output and prior proof version is
retained.
\end{remark}
